% GENERATED FILE. Edit canonical lessons, then run npm run generate:books.
% canonical-source-hash: fnv1a64:eff4a5bd39d9f0dd
% canonical-lessons: AR-C440-nafura, AR-C440-tidhkar, AR-C440-istiraha, AR-C440-rafaqa, AR-C440-awsala

\chapter{Breaks and souvenirs}
\label{ch:ar-breaks-and-souvenirs}
\hlchaptermodality{440}

\begin{chapteropening}
\textbf{By the end of this chapter:} \emph{I can say a fountain, a souvenir, a rest stop, to accompany and to give a lift.}

Complete the last lesson of chapter 440: I can say a fountain, a souvenir, a rest stop, to accompany and to give a lift.
\end{chapteropening}

\section[\texorpdfstring{\ar{نافورة} (nāfūra)}{نافورة (nāfūra)}]{\ar{نافورة} (nāfūra) — a fountain}

\label{lesson:AR-C440-nafura}



\begin{quote}
Before the new one: say the Arabic for to explore, then the Arabic for to take along.
\end{quote}

\subsection*{You'll want to know: \ar{نافورة}}
\textbf{\ar{نافورة}} — \emph{nāfūra} — \textquotedblleft{}a fountain\textquotedblright{}.

\begin{quote}
\ar{لعب} \ar{الأطفال} \ar{حول} \ar{النافورة}.
\end{quote}

\begin{quote}
\emph{laʿiba al-aṭfāl ḥawla an-nāfūra.}
\end{quote}

\begin{quote}
The children played around the fountain.
\end{quote}

\begin{cousinweb}
\textbf{\ar{نافورة}} (\emph{nāfūra}) names a fountain; the origin of the word is not settled.
\end{cousinweb}

\begin{grammarlens}[title={What you've built}]
One more word for getting around.
\end{grammarlens}

\subsection*{Your turn}
\begin{itemize}
\raggedright
  \item \emph{Say it:} \emph{nāfūra}
  \item \emph{Say it:} \emph{nāfūra}, once more
  \item \emph{Say it:} \emph{laʿiba al-aṭfāl ḥawla an-nāfūra.} — The children played around the fountain.
  \item \emph{Recall:} say the Arabic for to explore, then the Arabic for to take along, then say \emph{nāfūra} again
\end{itemize}

\begin{tcolorbox}[breakable,colback=teal!4,colframe=teal!35!black,title={Before you move on}]
What does \ar{نافورة} mean? (\textquotedblleft{}A fountain\textquotedblright{}.) Say it once more.
\end{tcolorbox}
\section[\texorpdfstring{\ar{تذكار} (tidhkār)}{تذكار (tidhkār)}]{\ar{تذكار} (tidhkār) — a souvenir; a keepsake}

\label{lesson:AR-C440-tidhkar}



\begin{quote}
Before the new one: say the Arabic for to take along, then the Arabic for a fountain.
\end{quote}

\subsection*{You'll want to know: \ar{تذكار}}
\textbf{\ar{تذكار}} — \emph{tidhkār} — \textquotedblleft{}a souvenir; a keepsake\textquotedblright{}.

\begin{quote}
\ar{اشتريت} \ar{تذكارا} \ar{صغيرا} \ar{لأختي}.
\end{quote}

\begin{quote}
\emph{ishtaraytu tidhkāran ṣaghīran li-ukhtī.}
\end{quote}

\begin{quote}
I bought a small souvenir for my sister.
\end{quote}

\begin{cousinweb}
\textbf{\ar{تذكار}} (\emph{tidhkār}) is built on the native Arabic root \emph{dh-k-r}, \textquotedblleft{}to remember\textquotedblright{} (Lane). A \emph{tidhkār} is a thing that makes one remember.
\end{cousinweb}

\begin{grammarlens}[title={What you've built}]
One more word for getting around.
\end{grammarlens}

\subsection*{Your turn}
\begin{itemize}
\raggedright
  \item \emph{Say it:} \emph{tidhkār}
  \item \emph{Say it:} \emph{tidhkār}, once more
  \item \emph{Say it:} \emph{ishtaraytu tidhkāran ṣaghīran li-ukhtī.} — I bought a small souvenir for my sister.
  \item \emph{Recall:} say the Arabic for to take along, then the Arabic for a fountain, then say \emph{tidhkār} again
\end{itemize}

\begin{tcolorbox}[breakable,colback=teal!4,colframe=teal!35!black,title={Before you move on}]
What does \ar{تذكار} mean? (\textquotedblleft{}A souvenir; a keepsake\textquotedblright{}.) Say it once more.
\end{tcolorbox}
\section[\texorpdfstring{\ar{استراحة} (istirāḥa)}{استراحة (istirāḥa)}]{\ar{استراحة} (istirāḥa) — a rest stop; a break}

\label{lesson:AR-C440-istiraha}



\begin{quote}
Before the new one: say the Arabic for a fountain, then the Arabic for a souvenir.
\end{quote}

\subsection*{You'll want to know: \ar{استراحة}}
\textbf{\ar{استراحة}} — \emph{istirāḥa} — \textquotedblleft{}a rest stop; a break\textquotedblright{}.

\begin{quote}
\ar{وقفت} \ar{الحافلة} \ar{عند} \ar{استراحة} \ar{على} \ar{الطريق}.
\end{quote}

\begin{quote}
\emph{waqafat al-ḥāfila ʿinda istirāḥa ʿalā aṭ-ṭarīq.}
\end{quote}

\begin{quote}
The bus stopped at a rest stop on the road.
\end{quote}

\begin{cousinweb}
\textbf{\ar{استراحة}} (\emph{istirāḥa}) is built on the native Arabic root \emph{r-w-ḥ}, \textquotedblleft{}to be at ease; breath, wind\textquotedblright{} (Lane). It is the verbal noun of Form X \emph{istarāḥa}, \textquotedblleft{}to rest\textquotedblright{}.
\end{cousinweb}

\begin{grammarlens}[title={What you've built}]
One more word for getting around.
\end{grammarlens}

\subsection*{Your turn}
\begin{itemize}
\raggedright
  \item \emph{Say it:} \emph{istirāḥa}
  \item \emph{Say it:} \emph{istirāḥa}, once more
  \item \emph{Say it:} \emph{waqafat al-ḥāfila ʿinda istirāḥa ʿalā aṭ-ṭarīq.} — The bus stopped at a rest stop on the road.
  \item \emph{Recall:} say the Arabic for a fountain, then the Arabic for a souvenir, then say \emph{istirāḥa} again
\end{itemize}

\begin{tcolorbox}[breakable,colback=teal!4,colframe=teal!35!black,title={Before you move on}]
What does \ar{استراحة} mean? (\textquotedblleft{}A rest stop; a break\textquotedblright{}.) Say it once more.
\end{tcolorbox}
\section[\texorpdfstring{\ar{رافق} (rāfaqa)}{رافق (rāfaqa)}]{\ar{رافق} (rāfaqa) — to accompany}

\label{lesson:AR-C440-rafaqa}



\begin{quote}
Before the new one: say the Arabic for a souvenir, then the Arabic for a rest stop.
\end{quote}

\subsection*{You'll want to know: \ar{رافق}}
\textbf{\ar{رافق}} — \emph{rāfaqa} — \textquotedblleft{}to accompany\textquotedblright{}.

\begin{quote}
\ar{رافقتنا} \ar{جدتي} \ar{في} \ar{الرحلة}.
\end{quote}

\begin{quote}
\emph{rāfaqatnā jaddatī fī ar-riḥla.}
\end{quote}

\begin{quote}
My grandmother came with us on the trip.
\end{quote}

Present: \ar{يرافق}, \emph{yurāfiq}.

\begin{cousinweb}
\textbf{\ar{رافق}} (\emph{rāfaqa}) is built on the native Arabic root \emph{r-f-q}, \textquotedblleft{}to be gentle; a companion\textquotedblright{} (Lane). It is Form III, \textquotedblleft{}to be a companion to\textquotedblright{}; a \emph{rafīq} is a companion.
\end{cousinweb}

\begin{grammarlens}[title={What you've built}]
One more word for getting around.
\end{grammarlens}

\subsection*{Your turn}
\begin{itemize}
\raggedright
  \item \emph{Say it:} \emph{rāfaqa}
  \item \emph{Say it:} \emph{rāfaqa}, once more
  \item \emph{Say it:} \emph{rāfaqatnā jaddatī fī ar-riḥla.} — My grandmother came with us on the trip.
  \item \emph{Recall:} say the Arabic for a souvenir, then the Arabic for a rest stop, then say \emph{rāfaqa} again
\end{itemize}

\begin{tcolorbox}[breakable,colback=teal!4,colframe=teal!35!black,title={Before you move on}]
What does \ar{رافق} mean? (\textquotedblleft{}To accompany\textquotedblright{}.) Say it once more.
\end{tcolorbox}
\section[\texorpdfstring{\ar{أوصل} (awṣala)}{أوصل (awṣala)}]{\ar{أوصل} (awṣala) — to give a lift; to deliver}

\label{lesson:AR-C440-awsala}



\begin{quote}
Before the new one: say the Arabic for a rest stop, then the Arabic for to accompany.
\end{quote}

\subsection*{You'll want to know: \ar{أوصل}}
\textbf{\ar{أوصل}} — \emph{awṣala} — \textquotedblleft{}to give a lift; to deliver\textquotedblright{}.

\begin{quote}
\ar{أوصلني} \ar{صديقي} \ar{إلى} \ar{المحطة}.
\end{quote}

\begin{quote}
\emph{awṣalanī ṣadīqī ilā al-maḥaṭṭa.}
\end{quote}

\begin{quote}
My friend gave me a lift to the station.
\end{quote}

Present: \ar{يوصل}, \emph{yūṣil}.

\begin{cousinweb}
\textbf{\ar{أوصل}} (\emph{awṣala}) is built on the native Arabic root \emph{w-ṣ-l}, \textquotedblleft{}to arrive, to reach\textquotedblright{} (Lane). It is Form IV, \textquotedblleft{}to make arrive\textquotedblright{}.
\end{cousinweb}

\begin{grammarlens}[title={What you've built}]
One more word for getting around.
\end{grammarlens}

\subsection*{Your turn}
\begin{itemize}
\raggedright
  \item \emph{Say it:} \emph{awṣala}
  \item \emph{Say it:} \emph{awṣala}, once more
  \item \emph{Say it:} \emph{awṣalanī ṣadīqī ilā al-maḥaṭṭa.} — My friend gave me a lift to the station.
  \item \emph{Recall:} say the Arabic for a rest stop, then the Arabic for to accompany, then say \emph{awṣala} again
\end{itemize}

\begin{tcolorbox}[breakable,colback=teal!4,colframe=teal!35!black,title={Before you move on}]
What does \ar{أوصل} mean? (\textquotedblleft{}To give a lift; to deliver\textquotedblright{}.) Say it once more.
\end{tcolorbox}
